% k0x_fill_skew: the inputs (P1)--(P3), the difference-kernel bound of lem:a07:K3 and \BALint\ on a lattice torus fitted to the blocks
% (input from sec/a07_k01steps_c.tex after rem:a07:dom). This is the transfer consumed by the rotation part on the tori of rot:hyp:F.

\subsection{The difference-kernel bound on a lattice torus fitted to the blocks}\label{ssec:a07:torus}

The rotation part applies King's \cite[Thm.~3.1]{K1}, and through it \BALint, on the tori $T_1=\DJ$ and $T_2=R^{-1}\DJ$ of \cref{rot:hyp:F}, each carrying both lattices $\Lk$ and $R^{-1}\Lk$ of \cref{rot:def:tori}(b). In the frame of the lattice $R^{-1}\Lk$ the torus $\DJ$ is the quotient of that lattice's own axis lattice by $R(50MJ\,\Z^3)$, and in the frame of $\Lk$ the torus $R^{-1}\DJ$ is the quotient by $R^{-1}(50MJ\,\Z^3)$: lattices of translations which are not rectangular in those frames. The statements of this appendix were written for subsets of $\eta\Z^3$. This subsection records that their proofs use the global geometry only through the two properties of \cref{def:a07:btorus}, which the tori of the rotation part have, so that \cref{lem:a07:CT,lem:a07:K3} and \cref{thm:a07:dom} hold on them with constants that do not depend on the torus.

\begin{definition}[a lattice torus fitted to the blocks]\label{def:a07:btorus}
In scale-$k$ units let $T=\eta\Z^3/\Gamma$, where $\Gamma\subset\Z^3$ is a lattice of rank $3$ (so $T$ carries the unit blocks and, since $\Gamma\subset\Z^3$, the blocks of every scale $j\le k$, of side $l_j=5^{j-k}$). Let $M_b\in\N$ be the side of a large block. $T$ is \emph{fitted to the blocks} if $\Gamma\subset M_b\Z^3$ and its sup-systole $\mathrm{sys}_\infty(T)\defeq\min\{|\gamma|_\infty:\gamma\in\Gamma\setminus\{0\}\}$ is at least $3M_b$. Distances on $T$ are the quotient sup-distances $\dist_T(x,x')\defeq\min_{\gamma\in\Gamma}|x-x'-\gamma|_\infty$; $\Omega\subset T$ denotes a union of large blocks, $T$ itself included, and ``regular on $\Omega$'' is \cite[(2.23)]{Ba1} read on the bonds of $\Omega$.
\end{definition}
If $T$ is fitted to the blocks in scale-$k$ units, so is its image in scale-$j$ units for every $j\le k$: the lattice becomes $l_j^{-1}\Gamma\subset l_j^{-1}M_b\Z^3\subset M_b\Z^3$ (as $l_j^{-1}=5^{k-j}\in\N$) and the systole grows by the factor $l_j^{-1}$. For the tori of the rotation part, in the frame of either lattice, the lattice of translations is $50MJ\,\Z^3$ or $R^{\pm1}(50MJ\,\Z^3)\subset10MJ\,\Z^3$ in physical units (\cref{rot:lem:geom}(b)). At Ba{\l}aban's step $k$ on the lattice of spacing $\eps_K=5^{-K}$ the scale-$k$ unit is $t=5^k\eps_K=5^{k-K}\le1$, so in scale-$k$ units these lattices lie in $10MJ\,5^{K-k}\,\Z^3$ and the systole is at least $10MJ\,5^{K-k}$. They are therefore fitted to the blocks at every step whenever the large-block side $M_b$ divides $10M$ and $3M_b\le10MJ$; in King's flow, the one run on these tori through \cite[Thm.~3.1]{K1}, the large block is King's and $M_b=M$.

\begin{lemma}[the printed inputs on a torus fitted to the blocks]\label{lem:a07:Ptorus}
Let $T$ be fitted to the blocks, $\Omega\subset T$ a union of large blocks and $A$ regular on $\Omega$. Then the bounds \textup{(P1)} (\cite[Prop.~2.1, (2.24)--(2.26)]{Ba1}, \cite[(1.9)--(1.12)]{Ba4}), \textup{(P2)} (\cite[Props.~2.2--2.3]{Ba1}, \cite[(1.13)--(1.20)]{Ba4}) and the positivity \cite[(1.8)]{Ba4} hold for the operators $G_k(\Omega,A)$, $\Delta^{(k)}(\Omega,A)$, $C^{(k)}_\Lambda(\Omega,A)$ on $T$, with the distances of \cref{def:a07:btorus}, with constants $c_0$, $\delta_0$, $R_0$, $c_2$, $\theta_0$, $\gamma$ that depend on $d$, $a$, $M_b$ and $L$ ($c_0$ also on $\alpha$) and not on $\Gamma$, $k$, $\Omega$, $\Lambda$ or $A$. For $\Omega=T$ the restriction to points at distance $\ge R_0$ from $\Omega^c$ is void. The same holds in scale-$j$ units for every $j\le k$.
\end{lemma}
\begin{proof}
Ba{\l}aban's proofs \cite[Sect.~2, pp.~574--586]{Ba4} are written for a union $\Omega$ of large blocks in $\eta\Z^d$ (``the torus identified with a rectangular parallelepiped''). They use the geometry of the ambient lattice in exactly the following places, each of which we verify on $T$.

\emph{(1) The covering cubes.} Let $v_j$ run through the vertices of the large-block grid of $T$ and $C_j\defeq v_j+[-M_b,M_b)^3$, the union of the $2^3$ large blocks with vertex $v_j$ (\cite[p.~578]{Ba4}: ``a cube with center $Mj$ and size equal to $2M$''). Since $2M_b<\mathrm{sys}_\infty(T)$, the projection $\eta\Z^3\to T$ is injective on each lift of $C_j$; it preserves bonds, unit blocks and the contours $\Gamma^{(k)}_{y,x}$ of the averaging operators, which lie inside blocks \cite[(1.4)]{Ba4}. Hence $\Pi_j\defeq\Omega\cap C_j$ is a union of at most $2^3$ large blocks in an embedded cube, the constant gauge $\lambda_0(x)=A_0\cdot(x-c_j)$ is defined on $C_j$, and \cite[Lemmas~2.1, 2.2, 2.4]{Ba4}, which are statements about $G_k(\Pi_j,A_j)$ on a parallelepiped of $\eta\Z^d$ with a local field $A_j$ built from $A|_{C_j}$, apply with their constants, which depend on $d$, $M_b$, $\alpha$ and the bare $a$ only.

\emph{(2) The partition of unity.} Let $h\in C^\infty_c((-\frac23,\frac23))$ with $h=1$ on $[-\frac13,\frac13]$ and $\sum_{n\in\Z}h(\cdot-n)^2=1$, and $h_j(x)\defeq\prod_\mu h((x_\mu-(v_j)_\mu)/M_b)$ \cite[p.~575]{Ba4}. Because $\Gamma\subset M_b\Z^3$, the family $\{h_j\}$ is $\Gamma$-periodic, descends to $T$, and $\sum_jh_j^2=1$ on $T$ with $\supp h_j\subset C_j$; so \cite[(2.2)--(2.9)]{Ba4} hold on $T$ as on $\eta\Z^d$ (if $\Omega=T$, then $\partial\Omega=\emptyset$ and (2.6) is trivial).

\emph{(3) The counts.} Each $j$ has at most $3^3$ neighbours and each point lies in at most $2^3$ of the supports, so the remainder of the parametrix has $\norm R_{2,2}\le c_dO(M_b^{-1})<1$ as on \cite[p.~577]{Ba4}, and the number $2^d(3^d)^{n-1}2^d$ of walks of length $n$ \cite[p.~579]{Ba4} is unchanged.

\emph{(4) The metric.} Consecutive cubes of a walk have centres at quotient distance $\le M_b$, so a walk of length $n$ from a cube meeting $\supp f$ to a cube containing $x$ has $n\ge\dist_T(x,\supp f)/M_b-2$; the geometric series \cite[(2.22)]{Ba4} then gives $\delta_0$, $c_0$ as on $\eta\Z^d$. The definition \cite[(2.19)]{Ba4} of $R_0$ and the implication ``$\dist_T(\{x,x'\},\Omega^c)\ge R_0$ implies that the first $n_0$ cubes of the walk lie in $\Omega$'' use only this. The cancellation of the interior-cube terms in $G_k(\Omega,A)-G_k(\Omega_0,A)$, which gives the two-sided factor of \cite[(2.26)]{Ba1}, uses (1) and (2). If $\Omega=T$, every $\Pi_j=C_j$ is a full cube and no restriction on $x$, $x'$ is needed \cite[p.~611]{Ba1}.

\emph{(5) Positivity.} For \cite[(1.8)]{Ba4}, $P_k(A)$ is block-diagonal, so dropping the bonds between unit blocks bounds the form from below by a sum over the unit blocks $\Delta$ of $\langle\phi,(-\Delta^N_{A,\Delta}+a_kP_k(A))\phi\rangle_\Delta$; on each block the gauge of (1) and the expansion in the field reduce the bound to $A=0$ \cite[p.~580]{Ba4}, where it is the Poincar\'e inequality on a unit block, which embeds in $T$. So $\gamma$ is Ba{\l}aban's.

The block operators of (P2) are built from $G_k(\Omega,A)$ and the averaging operators, and the proofs of \cite[Props.~2.2--2.3]{Ba1}, \cite[Sect.~3]{Ba4} use the same cubes, partition and walk and no further property of the ambient lattice; the restricted covariances $C^{(k)}_\Lambda$ for a union $\Lambda$ of large blocks are handled by the same walk on $\Lambda$. Nothing else in \cite[Sects.~2--3]{Ba4} refers to the global geometry. The last sentence of the lemma follows because the rescaled torus is fitted to the blocks (the remark after \cref{def:a07:btorus}) and the statements are the same at every scale.
\end{proof}

\begin{lemma}[\cref{lem:a07:CT,lem:a07:K3} and \BALint\ on a torus fitted to the blocks]\label{lem:a07:K3torus}
Let $T$ be fitted to the blocks. Then:
\begin{enumerate}[label=\textup{(\alph*)},nosep,leftmargin=*]
\item \cref{lem:a07:CT} holds for subsets of $T$ with the quotient metric, with the same constants.
\item \cref{lem:a07:K3} holds for unions $\Omega_2\subset\Omega$ of large blocks of $T$ and $A$ regular on $\Omega$, with $\dist_T$ in place of $\dist$ and with $C_{\dG}$, $\delta_{\dG}$ depending on $d$, $a$, $M_b$, $L$ and $\alpha$ only.
\item \cref{thm:a07:dom} holds on $T$, with the propagator domains $\OmA\supset\OmC$ and the background $\Bt$ of Ba{\l}aban's stage on $T$ (\cite[Sect.~2]{Ba2} run on $T$, as in \cref{thm:a06:BALA}) and with the renormalisation bound \cite[Prop.~1, (1.33)]{Ba3} on $T$; its constants $C_\kappa$, $t_0(\kappa)$ are those of the statement on $\eta\Z^3$ and do not depend on $\Gamma$.
\end{enumerate}
In particular all three hold on the tori $T_1$, $T_2$ of \cref{rot:hyp:F}, in the frame of either lattice.
\end{lemma}
\begin{proof}
(a) The proof of \cref{lem:a07:CT} uses the weight $\rho=\dist(\cdot,\supp f')$ only through two properties: $\rho$ is $1$-Lipschitz along bonds (the ends of a bond are at distance $\eta$), and each site is the final point of $d$ bonds. Both hold on $T$ with $\rho\defeq\dist_T(\cdot,\supp f')$, since $\dist_T(\cdot,S)$ is $1$-Lipschitz for the quotient metric and the bond structure is that of $\eta\Z^3$ locally. The lower bound $T\ge\gamma$ it needs is \cite[(1.8)]{Ba4} on $T$ (\cref{lem:a07:Ptorus}). The constants are unchanged.

(b) The proof of \cref{lem:a07:K3} consists of the slice decomposition $G_k(\Omega,A)=G^{[k]}_1+\sum_{j=1}^{k-1}S_j$, an operator identity which needs only that the blocks of every scale $j\le k$ partition $\Omega$ (true on $T$: a union of large blocks is a union of unit blocks of every scale $j\le k$, and these descend to $T$ because $\Gamma\subset M_b\Z^3\subset l_j\Z^3$); the rescaling to scale-$j$ units, under which $T$ stays fitted to the blocks; the slice bounds, which are (P1) at scale $j$ for the legs, (P2) at scale $j$ for the block covariances and (P3) at scale $1$, all available on $T$ with $\Gamma$-free constants by \cref{lem:a07:Ptorus} and (a); the sums over the block centres $\hat y,\hat y'$, which use \cref{lem:a07:G}(c) in the form $\sum_{\square}e^{-\delta\dist_T(\square,S)}\le C(\delta)|S|$, true on $T$ because the projection of the cubes of $\eta\Z^3$ onto those of $T$ is onto, so the sum on $T$ is at most the sum on $\eta\Z^3$ over a lift of $S$; and the triangle inequality for the metric. Every step is therefore available on $T$ with the quotient metric, and the proof goes through word for word. The constants $C_{\dG}$, $\delta_{\dG}$ are built from those of (P1)--(P3), $L$ and $\alpha$.

(c) In the proof of \cref{thm:a07:dom}: \cref{lem:a07:M,lem:a07:R,lem:a07:S} are statements about the expansion and regrouping of graph sums and about the sets of graphs forming a class; they do not refer to the geometry. \cref{lem:a07:C51} uses only that the three operators depend on the field through the bonds of their domain and on contours inside blocks. \cref{lem:a07:G}(a) uses the construction of the layers $\Lambda^{(k)}_i$ by sup-neighbourhoods of unions of large blocks and sup-metric geodesics, which exist in the quotient metric (a shortest path lifts to $\eta\Z^3$); (b) is a scalar inequality; (c) holds on $T$ as shown in (b) above, the number of cubes of $T$ at a given tree distance from a fixed cube being at most the number on $\eta\Z^3$. Steps 1--3 of the proof use (a) of \cref{lem:a07:G} and the definitions of the stage; Step 4 uses (b) above and (P1) on $T$; Step 5 uses \cite[Prop.~1, (1.33)]{Ba3} on $T$, which is Ba{\l}aban's renormalisation bound for a torus aligned with the blocks (\cref{thm:a06:BALA}); Steps 6 and 7 use \cref{lem:a07:G}(b), (c). The constants are those of the statement on $\eta\Z^3$.

For the tori of the rotation part, both lattices of translations are contained in $10MJ\,\Z^3$ in the frame in question (\cref{rot:lem:geom}(b)), so the tori are fitted to the blocks with $M_b=M$ by the remark after \cref{def:a07:btorus}.
\end{proof}

\begin{remark}\label{rem:a07:torus}
(i) The two conditions of \cref{def:a07:btorus} are what Ba{\l}aban uses of the rectangular geometry: that the large blocks tile the torus at every scale (the partition of unity and the block operators), and that a cube of two large blocks' side embeds (the local parametrix and the local gauge). A rectangular torus of sides that are multiples of $M_b$ and at least $3M_b$ satisfies both. (ii) The systole condition can be removed by passing to the covering torus $\eta\Z^3/3\Gamma$ and pulling back, at the price of the factor $3^3$ in the kernel bounds from the image sum over the deck group; it is not needed for the tori of the rotation part, whose systole is $10MJ$ at least.
\end{remark}
